%Einleitung

\begin{frame}
    \renewcommand{\figurename}{Figure}
    \fta{Introduction}
    
    \s{
       In the preceding modules, the fundamental physical principles and the basic components have been discussed.
        Furthermore, the essential technical behaviours in direct-current (DC) networks were clarified. In DC networks, electrical circuits are supplied with direct current, 
        whereas alternating-current (AC) networks operate with alternating current. DC flows continuously in one direction, while AC periodically reverses its direction. 
        The regularity of this reversal follows a specific frequency.
        For example, an ideal electrical resistor is not frequency-dependent. In contrast, components such as capacitors and inductors exhibit different complex impedances at different frequencies.
        To analyse these periodic quantities, the following topics are addressed:

        \begin{itemize}
            \item Pointer diagram
            \item Complex AC circuit analysis
            \item Effective value
            \item Apparent power, active power and reactive power
            \item Three-phase current
            \item Co-system, counter-system and zero system
        \end{itemize}

        \fu{
        \resizebox{0.7\textwidth}{!}{
        \input{Tikz/src/Einleitung/Zeigerdiagramm_Spannung_Strom}
        }
    }{{\bf Pointer diagram.} The quantities voltage and current are represented as vectors. \label{ZeigerBeispiel}}
    }

    \b{
        A sine wave has different voltage values over time. If the exact sine waves repeat themselves,
        it is a periodically recurring process. 
        Module 7, “Periodic quantities”, covers the following topics:\\

        \begin{minipage}[t]{0.45\textwidth}
            \begin{itemize}
                \vspace{\baselineskip}
                \item Zeigerdiagramm
                \item Komplexe Wechselstromrechnung
                \item Effektivwert
                \item Scheinleistung, Wirkleistung und Blindleistung
                \item Drehstrom
                \item Mitsystem, Gegensystem und Nullsystem
            \end{itemize}
        \end{minipage}
        \begin{minipage}[t]{0.54\textwidth}
        \fu{
            \resizebox{\textwidth}{!}{
            \begin{tikzpicture}
                \draw[very thin,color=gray] (-0.1,-0.1) grid (6.2,2.1);
                \draw[->] (-0.2,0) -- (6.4,0) node[right] {$\Re$};
                \draw[->] (0,-0.1) -- (0,2.2) node[above] {$\Im$};
                \draw[->, color=voltage] (0,0) coordinate (U) -- (6,0) coordinate (B) node[below] {$\underline{U}_\mathrm{1}$};
                \draw[->, color=red] (0,0.02) -- (3,0.02) node[below]	 {$\underline{I}_\mathrm{1}$};
                \draw[->, color=red] (3,0.02) -- (3,2.02) coordinate (D);
                \draw[color = red] (3,1)node[right]	 {$\underline{I}_\mathrm{c}$};
                \draw[->, color=red] (0,0.02) -- (3,2.02);
                \draw[color=red](2.3,1.9) node	 {$\underline{I}_\mathrm{2}$};
                \draw[->, color=voltage](6,0) -- (6,2) coordinate (C);
                \draw[color=voltage](6.2,1.3) node	 {$\underline{U}_\mathrm{2}$};
                \draw[->, color=voltage](0,0) -- (6,2);
                \draw[color=voltage](4.8,1.2) node	 {$\underline{U}_\mathrm{E}$};
                %Winkel
                \draw pic[draw, voltage, angle radius = 2cm, ->] {angle = B--U--C};
                \draw pic[draw, red, angle radius = 1.2cm, ->] {angle = B--U--D};
                \draw[red] (0.7,0.8) node{$\varphi_\mathrm{2}$};
                \draw[voltage] (2.32,0.2) node{$\varphi_\mathrm{1}$};
            \end{tikzpicture}
            }
        }{}
        \end{minipage}
    }
\end{frame}
\newpage